\section[Phân tích và thiết kế]{Phân tích và thiết kế}

\subsection{3.1 Khám phá dữ liệu}

\begin{frame}{Khám phá dữ liệu đo được 45 loại lỗi cần xử lý ở ODS}
\muc{3.1 · Khám phá dữ liệu}
\begin{columns}[T]\begin{column}{0.4\textwidth}
\begin{itemize}
  \item Cột thiếu, mốc thời gian sai thứ tự, tên thành phố lệch dấu, trùng lặp
\end{itemize}
\end{column}\begin{column}{0.58\textwidth}
\begin{center}\anh[0.74\textheight]{3-3-1_data_quality.png}\end{center}
\end{column}\end{columns}

\note{[S38] Không loại lỗi nào bằng cách xóa dữ liệu.}
\end{frame}

\begin{frame}{Tỷ lệ ô rỗng cho thấy cột nào cần điền unknown}
\muc{3.1 · Khám phá dữ liệu}
\begin{columns}[T]\begin{column}{0.4\textwidth}
\begin{itemize}
  \item Cột review và ngày giao có nhiều ô rỗng hợp lệ
\end{itemize}
\end{column}\begin{column}{0.58\textwidth}
\begin{center}\anh[0.74\textheight]{3-3-1_null_heatmap.png}\end{center}
\end{column}\end{columns}

\note{[S39] Ô rỗng của ngày giao đa phần là đơn chưa giao hoặc bị hủy.}
\end{frame}

\begin{frame}{Bảng toạ độ bẩn nhất: 93,5\% mã bưu chính có nhiều toạ độ, 26,2\% dòng trùng hoàn toàn}
\muc{3.1 · Khám phá dữ liệu}
\begin{columns}[T]\begin{column}{0.4\textwidth}
\begin{itemize}
  \item 7,3\% tên thành phố lệch dấu; 42 điểm ngoài lãnh thổ Brazil
\end{itemize}
\end{column}\begin{column}{0.58\textwidth}
\begin{center}\anh[0.74\textheight]{3-4-3_before_after_geolocation.png}\end{center}
\end{column}\end{columns}

\note{[S40] Giải pháp: lọc khung Brazil, lấy trung vị toạ độ và mode tên thành phố theo zip.}
\end{frame}

\begin{frame}{Phân bố giá, phí vận chuyển và lead time đều lệch phải}
\muc{3.1 · Khám phá dữ liệu}
\begin{columns}[T]\begin{column}{0.4\textwidth}
\begin{itemize}
  \item Nhiều giá trị lớn bất thường ở phí vận chuyển
\end{itemize}
\end{column}\begin{column}{0.58\textwidth}
\begin{center}\anh[0.74\textheight]{3-3-1_distributions_numeric.png}\end{center}
\end{column}\end{columns}

\note{[S41] Vì vậy dùng trung vị và Spearman khi cần.}
\end{frame}

\begin{frame}{Dữ liệu đơn dùng được là 01/2017–08/2018, hai đầu rất thưa}
\muc{3.1 · Khám phá dữ liệu}
\begin{columns}[T]\begin{column}{0.4\textwidth}
\begin{itemize}
  \item 2016 chỉ 329 đơn; 09–10/2018 chỉ 20 đơn
\end{itemize}
\end{column}\begin{column}{0.58\textwidth}
\begin{center}\anh[0.74\textheight]{3-3-1_orders_per_month.png}\end{center}
\end{column}\end{columns}

\note{[S42] Quyết định lọc thời gian đưa vào Power BI.}
\end{frame}

\begin{frame}{18/74 danh mục (24,3\%) tạo 80\% GMV, dẫn đầu health\symbol{95}beauty 9,3\%}
\muc{3.1 · Khám phá dữ liệu}
\begin{columns}[T]\begin{column}{0.4\textwidth}
\begin{itemize}
  \item Quy luật Pareto theo danh mục
\end{itemize}
\end{column}\begin{column}{0.58\textwidth}
\begin{center}\anh[0.74\textheight]{3-3-1_pareto_category.png}\end{center}
\end{column}\end{columns}

\note{[S43] Danh mục được gom thành 13 nhóm ngành hàng.}
\end{frame}

\begin{frame}{Seller thì phân tán hơn: 540 seller tạo 80\% GMV, 5 seller lớn nhất chỉ 7,7\%}
\muc{3.1 · Khám phá dữ liệu}
\begin{columns}[T]\begin{column}{0.4\textwidth}
\begin{itemize}
  \item Doanh thu không phụ thuộc vài seller
\end{itemize}
\end{column}\begin{column}{0.58\textwidth}
\begin{center}\anh[0.74\textheight]{3-3-1_pareto_seller.png}\end{center}
\end{column}\end{columns}

\note{[S44] Khác với danh mục, rủi ro tập trung vào seller thấp.}
\end{frame}

\begin{frame}{Giao trễ làm điểm review rơi từ 4,29 xuống 1,70 và 79,2\% review tiêu cực}
\muc{3.1 · Khám phá dữ liệu}
\begin{columns}[T]\begin{column}{0.4\textwidth}
\begin{itemize}
  \item Đơn trễ chiếm 6,5\% số review nhưng tạo 27,4\% review tiêu cực
\end{itemize}
\end{column}\begin{column}{0.58\textwidth}
\begin{center}\anh[0.74\textheight]{3-3-1_delay_vs_review.png}\end{center}
\end{column}\end{columns}

\note{[S45] Trung bình toàn sàn che mất nhóm nhỏ nhưng đau.}
\end{frame}

\begin{frame}{Trễ tập trung ở Bắc và Đông Bắc: Nordeste trễ 12,8\%, Norte lead time 22,6 ngày}
\muc{3.1 · Khám phá dữ liệu}
\begin{columns}[T]\begin{column}{0.4\textwidth}
\begin{itemize}
  \item Sudeste chỉ 6,1\% trễ, lead time 10,7 ngày
\end{itemize}
\end{column}\begin{column}{0.58\textwidth}
\begin{center}\anh[0.74\textheight]{3-3-1_late_rate_by_state.png}\end{center}
\end{column}\end{columns}

\note{[S46] Địa lý quyết định trải nghiệm giao hàng.}
\end{frame}

\begin{frame}{Chỉ 380/842 deal (45,1\%) có đơn: funnel và bán hàng chỉ giao nhau một phần}
\muc{3.1 · Khám phá dữ liệu}
\begin{columns}[T]\begin{column}{0.4\textwidth}
\begin{itemize}
  \item 380 seller tạo 674.005 BRL, 4,99\% GMV
  \item 380 là số seller có đơn ở mọi trạng thái; kho chỉ tính đơn không bị huỷ nên có 379 seller (45,0\%)
\end{itemize}
\end{column}\begin{column}{0.58\textwidth}
\begin{center}\anh[0.74\textheight]{3-3-1_funnel_seller_overlap.png}\end{center}
\end{column}\end{columns}

\note{[S47] Quyết định DB4 hiển thị cả nhóm "chốt nhưng chưa bán".}
\end{frame}

\begin{frame}{Sáu quyết định thiết kế đến từ khám phá dữ liệu}
\muc{3.1 · Khám phá dữ liệu}
\begin{itemize}
  \item SCD2 khách (0,26\% đổi zip); DB4 hai trang (cộng trang chi tiết); snapshot seller dày đặc; 13 nhóm ngành; zip thiếu toạ độ vẫn có dòng; mốc nạp 2017-12-31
\end{itemize}
\begin{center}\adjustbox{max width=\linewidth,max height=0.5\textheight}{\scriptsize\renewcommand{\arraystretch}{1.15}\begin{tabular}{lll}\toprule
\textbf{Quyết định} & \textbf{Nội dung} & \textbf{Số liệu đi kèm} \\\midrule
Câu a & \texttt{dim\_customer} dùng SCD2 & chỉ 0,26\% khách đổi zip \\
Câu b & DB4 hai trang & 380/842 deal có đơn \\
Câu c & snapshot seller dày đặc & 36.583 so với 16.441 dòng \\
Câu d & 13 nhóm ngành hàng & 73 danh mục \\
Câu e & zip thiếu toạ độ vẫn có dòng & 167 zip, không mất đơn \\
Câu f & mốc nạp đầu 2017-12-31 & 45.430 đơn, 11 đợt sau \\
\bottomrule\end{tabular}}\end{center}

\note{[S48] Mỗi quyết định có số liệu trong \texttt{handoff/design-questions.md}.}
\end{frame}

\begin{frame}{Mở rộng địa lý: hàm haversine SQL thuần tính khoảng cách seller $\rightarrow$ khách cho 112.650 dòng hàng}
\muc{3.1 · Khám phá dữ liệu}
\begin{itemize}
  \item 555 dòng (0,49\%) thiếu toạ độ vẫn được giữ lại với giá trị trống
\end{itemize}
\khungchen[0.3\textheight]{CHÈN HÌNH / BẢNG: đoạn SQL \texttt{dw.haversine\_\allowbreak{}km} (Phụ lục B)}

\note{[S107] São Paulo – Rio tính ra 360,9 km, gần với thực tế.}
\end{frame}

\begin{frame}{Khoảng cách tăng từ dưới 100 km lên trên 1.000 km làm phí vận chuyển gấp 2,7 lần và lead time gấp 2,9 lần}
\muc{3.1 · Khám phá dữ liệu}
\begin{itemize}
  \item Phí 11,75 $\rightarrow$ 31,65 BRL; lead time 6,5 $\rightarrow$ 19,0 ngày; trễ 4,5\% $\rightarrow$ 10,3\%
\end{itemize}
\begin{columns}[c]\begin{column}{0.48\textwidth}\anh[0.46\textheight]{3-5-2_distance_vs_freight.png}\end{column}\begin{column}{0.48\textwidth}\anh[0.46\textheight]{3-5-2_distance_vs_leadtime.png}\end{column}\end{columns}

\note{[S108] Tương quan Spearman 0,63 với phí và 0,54 với lead time; đó là tương quan, không phải nhân quả.}
\end{frame}

\begin{frame}{Haversine thuần đủ cho khoảng cách điểm–điểm; PostGIS cần khi truy vấn bán kính hoặc đa giác}
\muc{3.1 · Khám phá dữ liệu}
\begin{itemize}
  \item GiST chỉ mục không gian; haversine sai số tối đa 0,5\%
\end{itemize}
\khungchen[0.3\textheight]{CHÈN HÌNH / BẢNG: bảng so sánh (\texttt{report/chapters/3-1-kham-pha.tex})}

\note{[S109] Nhóm không cài PostGIS nên chọn giải pháp thuần SQL.}
\end{frame}

\begin{frame}{Hai phương pháp cùng đánh dấu 24.213 ngoại lai; nhóm giữ vì đó là ngoại lai hợp lệ}
\muc{3.1 · Khám phá dữ liệu}
\begin{columns}[T]\begin{column}{0.4\textwidth}
\begin{itemize}
  \item IQR đánh dấu 45.568, Isolation Forest (contamination 0,05, random\_state 42) 26.212
  \item Phí cao đi với hàng nặng (trung vị 5.700 g so với 700 g); 76,9\% đơn lead time ngoại lai là đơn trễ
  \item Chỉ 8 dòng cân nặng 0 g là lỗi dữ liệu
\end{itemize}
\end{column}\begin{column}{0.58\textwidth}
\begin{center}\anh[0.74\textheight]{3-1_outliers_boxplot.png}\end{center}
\end{column}\end{columns}

\note{[S133] Ngoại lai bằng IQR và iForest trên 5 biến.}
\end{frame}

\begin{frame}{Cramér's V lớn nhất 0,403 (giao trễ và nhóm review); tỉ số tương quan review theo cờ trễ 0,392}
\muc{3.1 · Khám phá dữ liệu}
\begin{columns}[T]\begin{column}{0.4\textwidth}
\begin{itemize}
  \item V hiệu chỉnh Bergsma cho cặp biến phân loại; $\eta$ cho biến phân loại và biến số
  \item Vùng khách -- vùng seller V = 0,072: gần như độc lập
  \item Spearman ngày trễ -- review chỉ $-0{,}18$: ảnh hưởng mang tính ngưỡng
\end{itemize}
\end{column}\begin{column}{0.58\textwidth}
\begin{center}\anh[0.74\textheight]{3-1_cramer_v_heatmap.png}\end{center}
\end{column}\end{columns}

\note{[S134] Công thức V, hiệu chỉnh Bergsma, tỉ số tương quan và thang đọc ở báo cáo mục 3.1.}
\end{frame}

\begin{frame}{STL: xu hướng mạnh 0,73, mùa vụ tuần 0,22; dưới 2 năm nên không kết luận mùa vụ theo năm}
\muc{3.1 · Khám phá dữ liệu}
\begin{columns}[T]\begin{column}{0.4\textwidth}
\begin{itemize}
  \item $F = \max(0, 1 - \mathrm{Var}(R)/\mathrm{Var}(S+R))$
  \item Black Friday 24/11/2017: 1.176 đơn, gấp 6,6 lần trung bình tháng 11
  \item Chu kỳ 12 tháng trên 20 điểm: $F_S = 0{,}80$ nhưng không đủ cơ sở
\end{itemize}
\end{column}\begin{column}{0.58\textwidth}
\begin{center}\anh[0.74\textheight]{3-1_stl_daily.png}\end{center}
\end{column}\end{columns}

\note{[S135] Phân rã STL số đơn theo ngày.}
\end{frame}

\subsection{3.2 Kiến trúc}

\begin{frame}{Kiến trúc sáu lớp từ hai hệ thống nguồn đến bốn dashboard}
\muc{3.2 · Kiến trúc}
\begin{columns}[T]\begin{column}{0.4\textwidth}
\begin{itemize}
  \item oltp/crm $\rightarrow$ staging $\rightarrow$ ODS $\rightarrow$ kho $\rightarrow$ data mart $\rightarrow$ Power BI; etl\_log giám sát
\end{itemize}
\end{column}\begin{column}{0.58\textwidth}
\begin{center}\anh[0.74\textheight]{3-4-1_pipeline.png}\end{center}
\end{column}\end{columns}

\note{[S49] Mỗi lớp có một việc và một cách kiểm tra riêng.}
\end{frame}

\begin{frame}{Tách lớp giúp đối soát: staging = ODS + dòng bị loại cho mọi bảng}
\muc{3.2 · Kiến trúc}
\begin{itemize}
  \item Chỗ nào hỏng biết ngay lớp nào
\end{itemize}
\khungchen[0.3\textheight]{CHÈN HÌNH / BẢNG: bảng đối soát (handoff P2)}

\note{[S50] Không phải lấy lại dữ liệu từ nguồn khi ODS lỗi.}
\end{frame}

\begin{frame}{Công nghệ: PostgreSQL 18, Python 3.11, SQL procedure, Power BI, Git}
\muc{3.2 · Kiến trúc}
\begin{itemize}
  \item Không Docker, không PostGIS; điều phối bằng một lệnh
\end{itemize}
\khungchen[0.3\textheight]{CHÈN HÌNH / BẢNG: logo công nghệ (chèn)}

\note{[S51] Chạy được trên laptop sinh viên 8 GB.}
\end{frame}

\begin{frame}{Kiến trúc theo lớp: sáu tầng từ nguồn đến BI, mỗi tầng có thành phần, công nghệ, số dòng}
\muc{3.2 · Kiến trúc}
\begin{columns}[T]\begin{column}{0.4\textwidth}
\begin{itemize}
  \item PostgreSQL, Python, Power BI, GitHub; số bảng và số dòng từng lớp
  \item Giám sát etl\_log 6 bảng; điều phối bằng một lệnh Python
\end{itemize}
\end{column}\begin{column}{0.58\textwidth}
\begin{center}\anh[0.74\textheight]{3-2_kien_truc_cac_lop.png}\end{center}
\end{column}\end{columns}

\note{[S136] Kiến trúc theo lớp (không vẽ luồng); luồng ở data flow và sơ đồ job.}
\end{frame}

\subsection{3.3 ETL}

\begin{frame}{ODS là lớp làm sạch trung gian: 18 bảng, 678.183 dòng, mọi lỗi có cờ hoặc lý do}
\muc{3.3 · ETL}
\begin{itemize}
  \item Zip thành 5 chữ số, thành phố bỏ dấu, \texttt{NULL} thành \texttt{unknown}, khử trùng review
\end{itemize}
\begin{center}\adjustbox{max width=\linewidth,max height=0.5\textheight}{\scriptsize\begin{tabular}{p{6.2cm}rp{5.5cm}}
\toprule
\textbf{Lỗi} & \textbf{Số dòng} & \textbf{Xử lý} \\
\midrule
Bàn giao vận chuyển trước khi duyệt thanh toán & 1.359 & Gắn cờ dq\_flags, giữ đơn; khoảng thời gian âm đặt NULL \\
Bàn giao vận chuyển trước ngày mua & 166 & Gắn cờ dq\_flags, giữ đơn; khoảng thời gian âm đặt NULL \\
Giao khách trước khi bàn giao vận chuyển & 23 & Gắn cờ dq\_flags, giữ đơn; khoảng thời gian âm đặt NULL \\
Đơn đã giao nhưng không có ngày duyệt & 14 & Gắn cờ dq\_flags, giữ đơn; khoảng thời gian âm đặt NULL \\
Trạng thái delivered nhưng không có ngày giao & 8 & Gắn cờ dq\_flags, giữ đơn; khoảng thời gian âm đặt NULL \\
Có ngày giao nhưng trạng thái chưa là delivered & 6 & Gắn cờ dq\_flags, giữ đơn; khoảng thời gian âm đặt NULL \\
Đã giao nhưng không có ngày bàn giao & 2 & Gắn cờ dq\_flags, giữ đơn; khoảng thời gian âm đặt NULL \\
Giá trị thanh toán = 0 & 9 & Gắn cờ, chuẩn hoá (số kỳ 0 → 1) \\
Hình thức thanh toán not\_defined & 3 & Gắn cờ, chuẩn hoá (số kỳ 0 → 1) \\
Số kỳ trả góp = 0 (đặt thành 1) & 2 & Gắn cờ, chuẩn hoá (số kỳ 0 → 1) \\
Doanh thu tháng khai báo = 0 & 797 & Gắn cờ, giữ nguyên giá trị \\
Chốt deal trước ngày liên hệ đầu & 1 & Gắn cờ, giữ nguyên giá trị \\
coordinates outside Brazil bounding box & 42 & Loại khỏi ODS, ghi etl\_log.rejected\_rows \\
\bottomrule
\end{tabular}}\end{center}

\note{[S52] 1.411 đơn (1,4\%) có ít nhất một cờ chất lượng.}
\end{frame}

\begin{frame}{Một lệnh chạy sáu bước ETL và ghi lại từng bước}
\muc{3.3 · ETL}
\begin{columns}[T]\begin{column}{0.4\textwidth}
\begin{itemize}
  \item extract $\rightarrow$ ODS $\rightarrow$ dim $\rightarrow$ fact $\rightarrow$ mart $\rightarrow$ tests $\rightarrow$ snapshot $\rightarrow$ export
\end{itemize}
\end{column}\begin{column}{0.58\textwidth}
\begin{center}\anh[0.74\textheight]{3-4-1_pipeline.png}\end{center}
\end{column}\end{columns}

\note{[S67] Mỗi bước ghi \texttt{job\_\allowbreak{}run}; lỗi thì dừng và đánh dấu failed.}
\end{frame}

\begin{frame}{Extract chỉ lấy dòng chưa có: đơn theo watermark, bảng nhỏ theo khoá}
\muc{3.3 · ETL}
\begin{columns}[T]\begin{column}{0.4\textwidth}
\begin{itemize}
  \item Chế độ full 52 giây; incremental khoảng 22 giây
\end{itemize}
\end{column}\begin{column}{0.58\textwidth}
\begin{center}\anh[0.74\textheight]{3-4-1_pipeline.png}\end{center}
\end{column}\end{columns}

\note{[S68] MQL và deal không dùng watermark ngày vì deal có thể đến trước MQL.}
\end{frame}

\begin{frame}{ODS xử lý lỗi theo ba cách: chuẩn hóa, gắn cờ hoặc loại có lý do}
\muc{3.3 · ETL}
\begin{itemize}
  \item 1.359 đơn bàn giao trước khi duyệt được gắn cờ, không xóa
\end{itemize}
\begin{center}\adjustbox{max width=\linewidth,max height=0.5\textheight}{\scriptsize\begin{tabular}{p{6.2cm}rp{5.5cm}}
\toprule
\textbf{Lỗi} & \textbf{Số dòng} & \textbf{Xử lý} \\
\midrule
Bàn giao vận chuyển trước khi duyệt thanh toán & 1.359 & Gắn cờ dq\_flags, giữ đơn; khoảng thời gian âm đặt NULL \\
Bàn giao vận chuyển trước ngày mua & 166 & Gắn cờ dq\_flags, giữ đơn; khoảng thời gian âm đặt NULL \\
Giao khách trước khi bàn giao vận chuyển & 23 & Gắn cờ dq\_flags, giữ đơn; khoảng thời gian âm đặt NULL \\
Đơn đã giao nhưng không có ngày duyệt & 14 & Gắn cờ dq\_flags, giữ đơn; khoảng thời gian âm đặt NULL \\
Trạng thái delivered nhưng không có ngày giao & 8 & Gắn cờ dq\_flags, giữ đơn; khoảng thời gian âm đặt NULL \\
Có ngày giao nhưng trạng thái chưa là delivered & 6 & Gắn cờ dq\_flags, giữ đơn; khoảng thời gian âm đặt NULL \\
Đã giao nhưng không có ngày bàn giao & 2 & Gắn cờ dq\_flags, giữ đơn; khoảng thời gian âm đặt NULL \\
Giá trị thanh toán = 0 & 9 & Gắn cờ, chuẩn hoá (số kỳ 0 → 1) \\
Hình thức thanh toán not\_defined & 3 & Gắn cờ, chuẩn hoá (số kỳ 0 → 1) \\
Số kỳ trả góp = 0 (đặt thành 1) & 2 & Gắn cờ, chuẩn hoá (số kỳ 0 → 1) \\
Doanh thu tháng khai báo = 0 & 797 & Gắn cờ, giữ nguyên giá trị \\
Chốt deal trước ngày liên hệ đầu & 1 & Gắn cờ, giữ nguyên giá trị \\
coordinates outside Brazil bounding box & 42 & Loại khỏi ODS, ghi etl\_log.rejected\_rows \\
\bottomrule
\end{tabular}}\end{center}

\note{[S69] Khoảng thời gian âm đặt NULL ở fact và được đếm lại.}
\end{frame}

\begin{frame}{Trước và sau ETL: bảng đơn hàng có thêm cờ lỗi và kiểu dữ liệu chuẩn}
\muc{3.3 · ETL}
\begin{columns}[T]\begin{column}{0.4\textwidth}
\begin{itemize}
  \item Ví dụ trước/sau
\end{itemize}
\end{column}\begin{column}{0.58\textwidth}
\begin{center}\anh[0.74\textheight]{3-4-3_before_after_orders.png}\end{center}
\end{column}\end{columns}

\note{[S70] Cùng 99.441 đơn nhưng thêm \texttt{dq\_\allowbreak{}flags}.}
\end{frame}

\begin{frame}{Trước và sau ETL: sản phẩm và seller được chuẩn hóa tên cột và thành phố}
\muc{3.3 · ETL}
\begin{itemize}
  \item \texttt{*\_\allowbreak{}lenght} $\rightarrow$ \texttt{*\_\allowbreak{}length}; thành phố bỏ dấu, chữ thường
\end{itemize}
\begin{columns}[c]\begin{column}{0.48\textwidth}\anh[0.46\textheight]{3-4-3_before_after_products.png}\end{column}\begin{column}{0.48\textwidth}\anh[0.46\textheight]{3-4-3_before_after_sellers.png}\end{column}\end{columns}

\note{[S71] Hai hình cùng một slide.}
\end{frame}

\begin{frame}{Trước và sau ETL: 1.000.163 dòng toạ độ gom thành 19.178 mã bưu chính}
\muc{3.3 · ETL}
\begin{columns}[T]\begin{column}{0.4\textwidth}
\begin{itemize}
  \item Trung vị toạ độ, mode thành phố
\end{itemize}
\end{column}\begin{column}{0.58\textwidth}
\begin{center}\anh[0.74\textheight]{3-4-3_before_after_geolocation.png}\end{center}
\end{column}\end{columns}

\note{[S72] 42 điểm ngoài Brazil được ghi vào bảng loại.}
\end{frame}

\begin{frame}{Trước và sau ETL: MQL và deal có giá trị trống thành unknown}
\muc{3.3 · ETL}
\begin{itemize}
  \item \texttt{origin} rỗng $\rightarrow$ \texttt{unknown}; doanh thu khai báo 0 gắn cờ
\end{itemize}
\begin{columns}[c]\begin{column}{0.48\textwidth}\anh[0.46\textheight]{3-4-3_before_after_mql.png}\end{column}\begin{column}{0.48\textwidth}\anh[0.46\textheight]{3-4-3_before_after_closed_deals.png}\end{column}\end{columns}

\note{[S73] Không đoán giá trị, chỉ đánh dấu.}
\end{frame}

\begin{frame}{Đối soát số dòng: nguồn = staging = ODS + loại, và tổng tiền khớp từng đồng}
\muc{3.3 · ETL}
\begin{itemize}
  \item $\Sigma$ price 13.591.643,70 BRL = file CSV
\end{itemize}
\khungchen[0.3\textheight]{CHÈN HÌNH / BẢNG: bảng đối soát (handoff P4)}

\note{[S74] Chênh lệch nếu có phải giải thích được bằng \texttt{rejected\_\allowbreak{}rows}.}
\end{frame}

\begin{frame}{SCD2 được nạp bằng window function: mỗi lần zip đổi mở một phiên bản mới}
\muc{3.3 · ETL}
\begin{itemize}
  \item \texttt{lag(zip)} so với đơn trước; \texttt{effective\_\allowbreak{}from} = ngày mua
\end{itemize}
\khungchen[0.3\textheight]{CHÈN HÌNH / BẢNG: đoạn SQL (Phụ lục B, \texttt{code/scd2\_\allowbreak{}dim\_\allowbreak{}customer.sql})}

\note{[S75] Khoá thay thế cũ không đổi khi chạy lại.}
\end{frame}

\begin{frame}{Nạp fact: tra khoá theo phiên bản có hiệu lực, khoảng thời gian âm thành NULL}
\muc{3.3 · ETL}
\begin{itemize}
  \item 1.359 đơn có \texttt{handling\_\allowbreak{}days} âm (bàn giao trước khi duyệt) được đặt NULL và đếm
\end{itemize}
\khungchen[0.3\textheight]{CHÈN HÌNH / BẢNG: bảng \texttt{fact\_\allowbreak{}load\_\allowbreak{}issue}}

\note{[S76] Tra không thấy thì dùng -1 và ghi nhận.}
\end{frame}

\begin{frame}{Nạp tăng dần mô phỏng: 11 đợt hằng tháng, tổng 286.144 bản ghi mới}
\muc{3.3 · ETL}
\begin{columns}[T]\begin{column}{0.4\textwidth}
\begin{itemize}
  \item \texttt{simulate\_\allowbreak{}increment} rồi \texttt{run\_\allowbreak{}pipeline --mode incremental}
\end{itemize}
\end{column}\begin{column}{0.58\textwidth}
\begin{center}\anh[0.74\textheight]{3-4-1_pipeline.png}\end{center}
\end{column}\end{columns}

\note{[S77] Đơn vào với trạng thái cuối; đó là giới hạn của mô phỏng.}
\end{frame}

\begin{frame}{Mỗi đợt chạy 20–26 giây, trong đó kiểm thử chiếm hơn một nửa}
\muc{3.3 · ETL}
\begin{itemize}
  \item Bảng 11 đợt: bản ghi mới, thời gian
\end{itemize}
\begin{center}\adjustbox{max width=\linewidth,max height=0.5\textheight}{\scriptsize\begin{tabular}{clrrrrrrrrr}
\toprule
\textbf{Đợt} & \textbf{Đến tháng} & \textbf{Bản ghi mới ở nguồn} & \textbf{ODS chèn} & \textbf{ODS loại} & \textbf{DW chèn / cập nhật} & \textbf{Giây} & \textbf{MQL} & \textbf{Deal} & \textbf{Seller có đơn} & \textbf{Đơn} \\
\midrule
1 & 2018-01 & 38.768 & 38.768 & 0 & 40.535 / 216 & 21,9 & 3.143 & 76 & 3 & 52.699 \\
2 & 2018-02 & 35.979 & 35.979 & 0 & 37.817 / 251 & 24,0 & 4.171 & 189 & 30 & 59.427 \\
3 & 2018-03 & 38.659 & 38.659 & 0 & 40.668 / 312 & 22,8 & 5.345 & 336 & 55 & 66.638 \\
4 & 2018-04 & 37.515 & 37.515 & 0 & 39.700 / 486 & 24,6 & 6.697 & 543 & 117 & 73.577 \\
5 & 2018-05 & 37.079 & 37.079 & 0 & 39.402 / 446 & 24,8 & 8.000 & 665 & 198 & 80.450 \\
6 & 2018-06 & 32.035 & 32.035 & 0 & 34.594 / 467 & 25,5 & 8.000 & 722 & 280 & 86.617 \\
7 & 2018-07 & 32.489 & 32.489 & 0 & 35.245 / 431 & 25,6 & 8.000 & 759 & 340 & 92.909 \\
8 & 2018-08 & 33.490 & 33.490 & 0 & 36.399 / 365 & 27,6 & 8.000 & 792 & 379 & 99.421 \\
9 & 2018-09 & 87 & 87 & 0 & 3.195 / 23 & 21,9 & 8.000 & 815 & 379 & 99.437 \\
10 & 2018-10 & 37 & 37 & 0 & 54 / 21 & 21,5 & 8.000 & 836 & 379 & 99.441 \\
11 & 2018-11 & 6 & 6 & 0 & 18 / 6 & 21,3 & 8.000 & 842 & 379 & 99.441 \\
\bottomrule
\end{tabular}}\end{center}

\note{[S78] Đợt 1–8 mang 99,95\% bản ghi mới; ba đợt cuối mang nốt 20 đơn và 50 deal.}
\end{frame}

\begin{frame}{Funnel lấp đầy dần: MQL 2.002 $\rightarrow$ 8.000, deal 3 $\rightarrow$ 842, seller có đơn 0 $\rightarrow$ 379}
\muc{3.3 · ETL}
\begin{itemize}
  \item Số liệu mỗi đợt từ \texttt{batch\_\allowbreak{}metrics}
\end{itemize}
\khungchen[0.3\textheight]{CHÈN HÌNH / BẢNG: biểu đồ đường từ \texttt{batch\_\allowbreak{}metrics} (bảng giám sát ở mục 3.3)}

\note{[S79] Đây là điểm trình bày của nạp tăng dần.}
\end{frame}

\begin{frame}{Kết quả incremental trùng nạp full: 16/16 bảng khớp số dòng, tổng tiền và nội dung}
\muc{3.3 · ETL}
\begin{itemize}
  \item So sánh bằng checksum trên khoá tự nhiên
\end{itemize}
\khungchen[0.3\textheight]{CHÈN HÌNH / BẢNG: bảng so sánh (\texttt{handoff/2026-11-05-P5.md})}

\note{[S80] Khoá thay thế khác nhau nên so sánh qua khoá tự nhiên.}
\end{frame}

\begin{frame}{Dòng mồ côi được thử lại ở đợt sau, nhưng cơ chế mới được chứng minh bằng kiểm thử}
\muc{3.3 · ETL}
\begin{itemize}
  \item Kịch bản deal \texttt{b91cf88…} trên bản sao trong giao dịch rồi hoàn tác
\end{itemize}
\khungchen[0.3\textheight]{CHÈN HÌNH / BẢNG: sơ đồ 3 bước (loại $\rightarrow$ cha đến $\rightarrow$ nạp lại)}

\note{[S81] Trong 11 đợt thật không có dòng mồ côi nào; nhóm nói thẳng điều này.}
\end{frame}

\begin{frame}{136 phép kiểm thử đạt 100\%, cộng 6 phép đối chiếu với khám phá dữ liệu}
\muc{3.3 · ETL}
\begin{itemize}
  \item Số dòng, tổng tiền, grain, khoá ngoại, SCD2, mart khớp fact
\end{itemize}
\begin{center}\adjustbox{max width=\linewidth,max height=0.5\textheight}{\scriptsize\renewcommand{\arraystretch}{1.15}\begin{tabular}{ll}\toprule
\textbf{Nhóm kiểm thử} & \textbf{Số phép} \\\midrule
ODS (số dòng, loại, cleaning) & 25 \\
Dimension (Unknown, SCD2, phủ khoá) & 42 \\
Fact (đối soát, grain, khoá ngoại, logic) & 49 \\
Data mart (khớp fact) & 17 \\
Thử lại dòng mồ côi (kịch bản) & 3 \\
\textbf{Tổng} & \textbf{136 (đạt 100\%)} \\
Đối chiếu khám phá dữ liệu (chỉ báo cáo) & 6/6 khớp \\
\bottomrule\end{tabular}}\end{center}

\note{[S82] Chạy sau mỗi lần pipeline, lưu \texttt{etl\_\allowbreak{}log.dq\_\allowbreak{}check}.}
\end{frame}

\begin{frame}{Nhật ký etl\symbol{95}log cho biết mỗi lô chạy gì, mất bao lâu, loại bao nhiêu dòng}
\muc{3.3 · ETL}
\begin{itemize}
  \item \texttt{job\_\allowbreak{}run}, \texttt{rejected\_\allowbreak{}rows}, \texttt{dq\_\allowbreak{}check}, \texttt{batch\_\allowbreak{}metrics}
\end{itemize}
\khungchen{CHÈN BẢNG: giám sát ETL (báo cáo mục 3.3)}

\note{[S83] Giám sát hệ thống nay là bảng ở mục 3.3 (không còn dashboard).}
\end{frame}

\begin{frame}{Dựng lại toàn bộ dự án từ đầu mất 443 giây và cho kết quả giống hệt}
\muc{3.3 · ETL}
\begin{itemize}
  \item Xóa hai database, chạy một lệnh, 16/16 bảng khớp
\end{itemize}
\begin{center}\adjustbox{max width=\linewidth,max height=0.5\textheight}{\scriptsize\begin{tabular}{p{6cm}rp{6cm}}
\toprule
\textbf{Bước} & \textbf{Giây} & \textbf{Kết quả} \\
\midrule
drop olist\_dw + olist\_full & 0.6 & dropped \\
setup\_db (schemas, DDL, procedures, views) & 1.8 & Setup done. \\
load\_source --initial (olist\_dw, to 2017-12-31) & 15.5 & Rejected rows logged in etl\_log.rejected\_rows: 0 (batch 1) \\
load\_source --full (olist\_full) & 24.0 & Rejected rows logged in etl\_log.rejected\_rows: 0 (batch 1) \\
run\_pipeline --mode full (initial state) & 55.8 & tests: 136/136 đạt, 0 lỗi; đối chiếu EDA 6/6 khớp; pipeline 54.7s \\
replay\_increments (11 x simulate + incremental) & 287.1 & 11 đợt hằng tháng (01/2018 → 11/2018); sau đợt cuối: 8,000 MQL, 842 deal, 379 seller có đơn bán được, 99,441 đơn, GMV 13,494,400.74 \\
run\_pipeline --mode incremental (tests + export\_powerbi) & 30.5 & tests: 136/136 đạt, 0 lỗi; đối chiếu EDA 6/6 khớp; pipeline 29.3s \\
\bottomrule
\end{tabular}}\end{center}

\note{[S110] Khả năng tái lập là bằng chứng độ tin cậy của pipeline.}
\end{frame}

\begin{frame}{Sáu giai đoạn ETL: lô đầu nạp 1.273.620 dòng vào staging, 42 dòng loại, 142 test đạt}
\muc{3.3 · ETL}
\begin{columns}[T]\begin{column}{0.4\textwidth}
\begin{itemize}
  \item Extract, seed và ODS, dimension và fact, mart, kiểm thử, xuất CSV
  \item Số dòng lấy từ etl\_log.job\_run (lô đầu và tổng các lô)
\end{itemize}
\end{column}\begin{column}{0.58\textwidth}
\begin{center}\anh[0.74\textheight]{3-3_etl_jobs.png}\end{center}
\end{column}\end{columns}

\note{[S137] Sơ đồ các job của run_pipeline sinh từ etl_log.}
\end{frame}

\begin{frame}{SQL, Python và Power Query cho cùng kết quả làm sạch: 8.011 thành 5.968 thành phố}
\muc{3.3 · ETL}
\begin{itemize}
  \item Chữ thường, bỏ dấu, zip đủ 5 ký tự; 73.440 dòng đổi
  \item Python chênh 2 dòng do ký tự ordinal; mã M thêm bảng thay thế để khớp SQL
  \item Mã M và hướng dẫn chụp before/after: powerbi/power-query-vi-du.md
\end{itemize}

\note{[S138] Nhóm cần chạy mã M trong Power BI và chụp ảnh before/after.}
\end{frame}

\begin{frame}{Giám sát ETL: 11 đợt tăng dần 21--28 giây, 136 test đạt 100\%, chỉ 42 dòng bị loại}
\muc{3.3 · ETL}
\begin{itemize}
  \item Dòng bị loại duy nhất là 42 điểm toạ độ ngoài Brazil; không có dòng mồ côi
  \item Funnel lấp đầy dần: MQL 2.002 $\rightarrow$ 8.000, deal 3 $\rightarrow$ 842, seller có đơn 0 $\rightarrow$ 379
  \item Bảng số dòng qua từng lớp, dòng bị loại, thời gian và test đạt từng đợt ở báo cáo mục 3.3 (trước đây là một dashboard)
\end{itemize}
\khungchen[0.3\textheight]{CHÈN BẢNG: các đợt nạp (tables/p8b\_etl\_batches)}

\note{[S105] Số liệu giám sát lấy từ etl_log.}
\end{frame}

\subsection{3.4 Chiều khái niệm}

\begin{frame}{Mười bảng chiều với số dòng từ 6 đến 96.339}
\muc{3.4 · Chiều khái niệm}
\begin{itemize}
  \item dim\_date 1.828; geography 19.178; customer 96.339; seller 3.558; product 32.952; ...
\end{itemize}
\begin{center}\adjustbox{max width=\linewidth,max height=0.5\textheight}{\scriptsize\begin{tabular}{llrc}
\toprule
\textbf{Chiều} & \textbf{Grain} & \textbf{Số dòng} & \textbf{SCD} \\
\midrule
dim\_date & 1 ngày (2016-01-01 → 2020-12-31) & 1.828 & -- \\
dim\_geography & 1 mã bưu chính (zip) & 19.178 & 1 \\
dim\_customer & 1 phiên bản địa chỉ của 1 khách thật & 96.339 & 2 \\
dim\_seller & 1 seller (chiều dùng chung) & 3.558 & 1 \\
dim\_product & 1 sản phẩm & 32.952 & 1 \\
dim\_order\_status & 1 trạng thái đơn & 9 & 1 \\
dim\_payment\_type & 1 hình thức thanh toán & 6 & 1 \\
dim\_lead\_origin & 1 giá trị origin của MQL & 11 & 1 \\
dim\_lead\_profile & 1 tổ hợp thuộc tính deal (junk dimension) & 503 & -- \\
dim\_sales\_rep & 1 nhân viên và vai trò (SDR/SR) & 55 & -- \\
\bottomrule
\end{tabular}}\end{center}

\note{[S57] Mỗi chiều có đúng một dòng Unknown, được kiểm thử.}
\end{frame}

\begin{frame}{Hệ thống chiều khái niệm: phân cấp ngày, địa lý, sản phẩm, seller, khách}
\muc{3.4 · Chiều khái niệm}
\begin{itemize}
  \item Ngày $\rightarrow$ tháng $\rightarrow$ quý $\rightarrow$ năm; zip $\rightarrow$ thành phố $\rightarrow$ bang $\rightarrow$ vùng
\end{itemize}
\khungchen[0.3\textheight]{CHÈN HÌNH / BẢNG: sơ đồ phân cấp (vẽ)}

\note{[S58] Phân cấp là cơ sở của roll-up và drill-down.}
\end{frame}

\begin{frame}{SCD2 giữ lịch sử địa chỉ: 96.096 khách thành 96.339 dòng, 236 khách đổi zip}
\muc{3.4 · Chiều khái niệm}
\begin{itemize}
  \item Đơn cũ vẫn trỏ về địa chỉ lúc mua
\end{itemize}
\begin{center}\adjustbox{max width=\linewidth,max height=0.5\textheight}{\scriptsize\renewcommand{\arraystretch}{1.15}\begin{tabular}{llllllll}\toprule
\textbf{customer\_key} & \textbf{khách thật} & \textbf{zip} & \textbf{thành phố} & \textbf{effective\_from} & \textbf{effective\_to} & \textbf{is\_current} & \textbf{version} \\\midrule
15081 & 570eb70f… & 31930 & belo horizonte & 1900-01-01 & 2017-11-08 & không & 1 \\
15082 & 570eb70f… & 31130 & belo horizonte & 2017-11-09 & 9999-12-31 & có & 2 \\
\bottomrule\end{tabular}}\end{center}

\note{[S59] Ví dụ khách \texttt{570eb70f…} đổi zip ngày 09/11/2017; đơn 24/10 thuộc phiên bản 1.}
\end{frame}

\begin{frame}{dim\symbol{95}seller là chiều dùng chung: nối bán hàng với funnel qua seller\symbol{95}id}
\muc{3.4 · Chiều khái niệm}
\begin{itemize}
  \item 462 seller chỉ có trong CRM vẫn có dòng, địa chỉ Unknown
\end{itemize}
\khungchen[0.3\textheight]{CHÈN HÌNH / BẢNG: bảng ví dụ seller (\texttt{huong-dan/01-pipeline-bang-loi.md})}

\note{[S60] \texttt{first\_\allowbreak{}sale\_\allowbreak{}date} được điền sau khi nạp fact.}
\end{frame}

\begin{frame}{Chiều rác và role-playing giúp fact gọn mà vẫn đủ thông tin}
\muc{3.4 · Chiều khái niệm}
\begin{itemize}
  \item \texttt{dim\_\allowbreak{}lead\_\allowbreak{}profile} gom 7 thuộc tính thành 503 tổ hợp; \texttt{dim\_\allowbreak{}date} đóng 5 vai ở fulfillment
\end{itemize}
\khungchen[0.3\textheight]{CHÈN HÌNH / BẢNG: bảng ví dụ hai kỹ thuật}

\note{[S61] Trong Power BI, vai trò ngày dùng quan hệ không hoạt động và \texttt{USERELATIONSHIP}.}
\end{frame}

\begin{frame}{Dòng Unknown -1 và khoá thay thế bảo đảm fact không bao giờ có khoá rỗng}
\muc{3.4 · Chiều khái niệm}
\begin{itemize}
  \item Mọi lỗi tra khoá được ghi \texttt{fact\_\allowbreak{}load\_\allowbreak{}issue}
\end{itemize}
\khungchen[0.3\textheight]{CHÈN HÌNH / BẢNG: bảng \texttt{dim\_\allowbreak{}seller} dòng -1}

\note{[S66] Nhờ đó ETL không làm rơi dòng khi chiều thiếu.}
\end{frame}

\begin{frame}{Phân cấp đúng khi tháng kèm năm, thành phố kèm bang: 303 tên thành phố trùng giữa các bang}
\muc{3.4 · Chiều khái niệm}
\begin{itemize}
  \item 12 số tháng nhưng 60 cặp (năm, tháng); 5.814 tên thành phố nhưng 6.190 cặp (bang, thành phố)
  \item Mỗi danh mục một nhóm ngành, mỗi origin một nhóm kênh (0 vi phạm)
  \item Mười chiều chia nhóm: dùng chung, đóng vai, SCD2, rác, danh mục nhỏ
\end{itemize}

\note{[S139] Bảng phân cấp từng chiều với giá trị ví dụ và số giá trị phân biệt ở báo cáo mục 3.4.}
\end{frame}

\subsection{3.5 Data model OLAP}

\begin{frame}{Bus matrix: sáu fact dùng chung các chiều ngày, seller, khách, địa lý}
\muc{3.5 · Data model OLAP}
\begin{itemize}
  \item Seller là chiều dùng chung của hai quy trình
\end{itemize}
\begin{center}\adjustbox{max width=\linewidth,max height=0.5\textheight}{\scriptsize\renewcommand{\arraystretch}{1.15}\begin{tabular}{lcccccccccc}\toprule
\textbf{Fact} & \textbf{Date} & \textbf{Customer} & \textbf{Seller} & \textbf{Product} & \textbf{Geo} & \textbf{Status} & \textbf{Payment} & \textbf{Origin} & \textbf{Profile} & \textbf{Rep} \\\midrule
fact\_order\_item & \ding{51} & \ding{51} & \ding{51} & \ding{51} & \ding{51} & \ding{51} &  &  &  &  \\
fact\_order\_fulfillment & \ding{51} & \ding{51} &  &  & \ding{51} & \ding{51} &  &  &  &  \\
fact\_payment & \ding{51} & \ding{51} &  &  &  &  & \ding{51} &  &  &  \\
fact\_review & \ding{51} & \ding{51} &  &  &  &  &  &  &  &  \\
fact\_lead\_funnel & \ding{51} &  & \ding{51} &  &  &  &  & \ding{51} & \ding{51} & \ding{51} \\
fact\_seller\_monthly & \ding{51} &  & \ding{51} &  &  &  &  &  &  &  \\
\bottomrule\end{tabular}}\end{center}

\note{[S53] Bus matrix cho thấy chiều nào được dùng lại ở đâu.}
\end{frame}

\begin{frame}{Mô hình constellation: bốn mart, sáu fact (cam), mười chiều (xanh)}
\muc{3.5 · Data model OLAP}
\begin{columns}[T]\begin{column}{0.4\textwidth}
\begin{itemize}
  \item Khung = data mart
\end{itemize}
\end{column}\begin{column}{0.58\textwidth}
\begin{center}\anh[0.74\textheight]{3-3-4_constellation.png}\end{center}
\end{column}\end{columns}

\note{[S54] Fact màu cam, chiều màu xanh đúng yêu cầu trình bày.}
\end{frame}

\begin{frame}{Data model logic: mỗi fact nối đúng các chiều theo bus matrix}
\muc{3.5 · Data model OLAP}
\begin{columns}[T]\begin{column}{0.4\textwidth}
\begin{itemize}
  \item Ví dụ \texttt{fact\_\allowbreak{}order\_\allowbreak{}item} nối 6 chiều với 2 vai địa lý
\end{itemize}
\end{column}\begin{column}{0.58\textwidth}
\begin{center}\anh[0.74\textheight]{3-3-4_erd_dw.png}\end{center}
\end{column}\end{columns}

\note{[S55] Địa lý có hai vai: khách và seller.}
\end{frame}

\begin{frame}{Data model vật lý: khoá chính theo grain, khoá ngoại có index, 112.650 dòng hàng}
\muc{3.5 · Data model OLAP}
\begin{columns}[T]\begin{column}{0.4\textwidth}
\begin{itemize}
  \item Mọi khoá ngoại NOT NULL, dòng Unknown -1
\end{itemize}
\end{column}\begin{column}{0.58\textwidth}
\begin{center}\anh[0.74\textheight]{3-3-4_erd_dw.png}\end{center}
\end{column}\end{columns}

\note{[S56] Tiền \texttt{numeric(12,2)}, khoá \texttt{<tên>\_\allowbreak{}key}, ngày \texttt{yyyymmdd}.}
\end{frame}

\begin{frame}{Sáu fact thuộc đủ ba loại của Kimball}
\muc{3.5 · Data model OLAP}
\begin{itemize}
  \item Transaction $\times$3, accumulating $\times$2, periodic $\times$1
\end{itemize}
\begin{center}\adjustbox{max width=\linewidth,max height=0.5\textheight}{\scriptsize\begin{tabular}{lllr}
\toprule
\textbf{Fact} & \textbf{Loại} & \textbf{Grain} & \textbf{Số dòng} \\
\midrule
fact\_order\_item & Transaction & 1 dòng hàng của 1 đơn & 112.650 \\
fact\_order\_fulfillment & Accumulating snapshot & 1 đơn (5 mốc thời gian) & 99.441 \\
fact\_payment & Transaction & 1 lần thanh toán & 103.886 \\
fact\_review & Transaction & 1 review của 1 đơn & 99.224 \\
fact\_lead\_funnel & Accumulating snapshot & 1 MQL & 8.000 \\
fact\_seller\_monthly & Periodic snapshot (dày đặc) & 1 seller × 1 tháng & 36.583 \\
\bottomrule
\end{tabular}}\end{center}

\note{[S62] Tổng 459.784 dòng fact (99.441 + 112.650 + 103.886 + 99.224 + 8.000 + 36.583).}
\end{frame}

\begin{frame}{fact\symbol{95}order\symbol{95}fulfillment là accumulating snapshot: một đơn một dòng, cập nhật theo mốc}
\muc{3.5 · Data model OLAP}
\begin{itemize}
  \item Mốc chưa xảy ra có khoá -1; thời gian từng chặng tính sẵn
\end{itemize}
\begin{center}\adjustbox{max width=\linewidth,max height=0.5\textheight}{\scriptsize\renewcommand{\arraystretch}{1.15}\begin{tabular}{ll}\toprule
\textbf{Cột của fact\_order\_fulfillment (đơn 19c36691…)} & \textbf{Giá trị} \\\midrule
approval\_hours & 36,61 \\
handling\_days & 0,71 \\
transit\_days & 12,84 \\
lead\_days & 15,07 \\
promised\_days & 21 \\
delay\_days & -6 (giao sớm 6 ngày) \\
order\_value & 43,00 \\
freight\_total & 45,30 \\
payment\_total & 88,30 \\
\bottomrule\end{tabular}}\end{center}

\note{[S63] Nếu dùng transaction fact phải ghép nhiều dòng mới tính được lead time.}
\end{frame}

\begin{frame}{fact\symbol{95}lead\symbol{95}funnel ghi lại hành trình MQL: ba mốc ngày và hai khoảng thời gian}
\muc{3.5 · Data model OLAP}
\begin{itemize}
  \item \texttt{days\_\allowbreak{}to\_\allowbreak{}close} TB 48,5; \texttt{days\_\allowbreak{}won\_\allowbreak{}to\_\allowbreak{}first\_\allowbreak{}sale} TB 51,9
\end{itemize}
\khungchen[0.3\textheight]{CHÈN HÌNH / BẢNG: dòng ví dụ MQL}

\note{[S64] Cũng là accumulating snapshot, 8.000 dòng.}
\end{frame}

\begin{frame}{fact\symbol{95}seller\symbol{95}monthly lưu dày đặc 36.583 dòng để thấy seller ngừng bán}
\muc{3.5 · Data model OLAP}
\begin{itemize}
  \item Dày đặc so với thưa: 36.583 và 16.441 dòng
\end{itemize}
\khungchen[0.3\textheight]{CHÈN HÌNH / BẢNG: bảng 4 tháng đầu của một seller}

\note{[S65] Tháng không có đơn vẫn có dòng với \texttt{orders = 0}.}
\end{frame}

\begin{frame}{Bốn data mart gồm 17 view, trong đó 4 là materialized view làm mới trong 2 giây}
\muc{3.5 · Data model OLAP}
\begin{itemize}
  \item Sales, fulfillment, cx, acquisition
\end{itemize}
\begin{center}\adjustbox{max width=\linewidth,max height=0.5\textheight}{\scriptsize\renewcommand{\arraystretch}{1.15}\begin{tabular}{lll}\toprule
\textbf{Mart} & \textbf{Số view} & \textbf{Dashboard} \\\midrule
mart\_sales & 6 & DB1, DB3 \\
mart\_fulfillment & 4 & DB2, DB3 \\
mart\_cx & 4 (RFM, cohort materialized) & DB2 \\
mart\_acquisition & 3 (scorecard, theo kênh materialized) & DB4 \\
\bottomrule\end{tabular}}\end{center}

\note{[S84] Mart được làm mới tự động ở bước mart của pipeline.}
\end{frame}

\begin{frame}{Roll-up và drill-down: GMV 2017 là 6,11 triệu, 2018 là 7,34 triệu; ngày Black Friday cao gấp 5,2 lần}
\muc{3.5 · Data model OLAP}
\begin{itemize}
  \item Quý 4/2017: 2,41 triệu BRL; 24/11/2017: 152.654 BRL
\end{itemize}
\begin{center}\adjustbox{max width=\linewidth,max height=0.5\textheight}{\scriptsize\begin{tabular}{llrr}
\toprule
\textbf{Năm} & \textbf{Quý} & \textbf{Số đơn} & \textbf{GMV (BRL)} \\
\midrule
2017 & 2017-Q1 & 5.122 & 733.398,94 \\
2017 & 2017-Q2 & 9.222 & 1.286.918,78 \\
2017 & 2017-Q3 & 12.445 & 1.681.949,00 \\
2017 & 2017-Q4 & 17.586 & 2.406.225,55 \\
2017 & (cả năm) & 44.375 & 6.108.492,27 \\
2018 & 2018-Q1 & 20.979 & 2.764.402,78 \\
2018 & 2018-Q2 & 19.897 & 2.849.730,26 \\
2018 & 2018-Q3 & 12.655 & 1.727.049,37 \\
2018 & (cả năm) & 53.531 & 7.341.182,41 \\
Tổng & (cả năm) & 97.906 & 13.449.674,68 \\
\bottomrule
\end{tabular}}\end{center}

\note{[S85] Roll-up dùng \texttt{ROLLUP}, drill-down thêm cấp vào \texttt{GROUP BY}.}
\end{frame}

\begin{frame}{Slice, dice và pivot: vùng Sudeste chiếm 64,3\% GMV năm 2017}
\muc{3.5 · Data model OLAP}
\begin{itemize}
  \item Dice Sudeste $\times$ Home \& Furniture $\times$ quý 4: SP 199.179 BRL
\end{itemize}
\begin{center}\adjustbox{max width=\linewidth,max height=0.5\textheight}{\scriptsize\begin{tabular}{lrrrr}
\toprule
\textbf{Vùng} & \textbf{2016} & \textbf{2017} & \textbf{2018} & \textbf{Tổng} \\
\midrule
Sudeste & 27.701 & 3.928.343 & 4.866.293 & 8.822.336 \\
Sul & 7.456 & 883.020 & 1.047.636 & 1.938.111 \\
Nordeste & 6.013 & 727.733 & 804.568 & 1.538.314 \\
Centro-Oeste & 2.356 & 401.380 & 458.090 & 861.825 \\
Norte & 1.200 & 168.018 & 164.596 & 333.814 \\
\bottomrule
\end{tabular}}\end{center}

\note{[S86] Pivot bằng tổng có điều kiện \texttt{FILTER}.}
\end{frame}

\begin{frame}{Chọn ROLAP vì dữ liệu nhỏ, kho luôn tươi và Power BI làm lớp ngữ nghĩa}
\muc{3.5 · Data model OLAP}
\begin{itemize}
  \item Khác MOLAP: không phải xử lý lại khối khi nạp dữ liệu
\end{itemize}
\khungchen[0.3\textheight]{CHÈN HÌNH / BẢNG: bảng ROLAP vs MOLAP (\texttt{handoff/olap-examples.md})}

\note{[S87] Đây là lý do lựa chọn đã nêu ở phần 1.}
\end{frame}

\begin{frame}{Star schema chỉ nhanh hơn OLTP 1,0–1,6 lần; tốc độ lớn đến từ materialized view và index}
\muc{3.5 · Data model OLAP}
\begin{columns}[T]\begin{column}{0.4\textwidth}
\begin{itemize}
  \item Chọn lọc một seller nhanh gấp 9 lần nhờ index; materialized view dưới 1 ms so với 150–750 ms
\end{itemize}
\end{column}\begin{column}{0.58\textwidth}
\begin{center}\anh[0.74\textheight]{3-5-1_benchmark.png}\end{center}
\end{column}\end{columns}

\note{[S88] Không phóng đại: giá trị chính của mô hình chiều là truy vấn đơn giản và số liệu nhất quán.}
\end{frame}

\begin{frame}{Data model logic: sáu quy trình nối với mười dimension và các cấp phân cấp}
\muc{3.5 · Data model OLAP}
\begin{itemize}
  \item Fact cam, dimension xanh; không ghi kiểu dữ liệu
\end{itemize}
\begin{center}\anh[0.62\textheight]{3-5_data_model_logic.png}\end{center}

\note{[S140] Mô hình logic.}
\end{frame}

\begin{frame}{Data model vật lý sinh từ cơ sở dữ liệu: mọi bảng dw với cột, kiểu, PK và FK}
\muc{3.5 · Data model OLAP}
\begin{itemize}
  \item 16 bảng, fact và dimension khác màu
\end{itemize}
\begin{center}\anh[0.62\textheight]{3-5_data_model_vat_ly.png}\end{center}

\note{[S141] Mô hình vật lý tổng thể.}
\end{frame}

\begin{frame}{Bốn data mart, mỗi mart gồm fact, dimension và 3--4 câu hỏi phân tích}
\muc{3.5 · Data model OLAP}
\begin{itemize}
  \item mart\_sales, mart\_fulfillment, mart\_cx, mart\_acquisition
\end{itemize}
\begin{center}\anh[0.62\textheight]{3-5_data_mart.png}\end{center}

\note{[S142] Data mart.}
\end{frame}

